\documentclass[12pt]{article}

\usepackage{amssymb, amsmath, fullpage, amsthm, fancybox}
\usepackage[demo]{graphicx}
\newtheorem*{theorem}{Theorem}
\newtheorem*{lemma}{Lemma}
\newtheorem*{proposition}{Proposition}
\newtheorem*{corollary}{Corollary}

%\newenvironment{proof}[1][Proof]{\begin{trivlist}
%\item[\hskip \labelsep {\bfseries #1}]}{\end{trivlist}}
\newenvironment{definition}[1][Definition]{\begin{trivlist}
\item[\hskip \labelsep {\bfseries #1}]}{\end{trivlist}}
\newenvironment{example}[1][Example]{\begin{trivlist}
\item[\hskip \labelsep {\bfseries #1}]}{\end{trivlist}}
\newenvironment{remark}[1][Remark]{\begin{trivlist}
\item[\hskip \labelsep {\bfseries #1}]}{\end{trivlist}}
\newcommand{\dfdx}{\frac{\partial f}{\partial x}}
\newcommand{\QED}{\nobreak \ifvmode \relax \else
     \ifdim\lastskip<1.5em \hskip-\lastskip
    \hskip1.5em plus0em minus0.5em \fi \nobreak
   \vrule height0.75em width0.5em depth0.25em\fi}
\renewcommand{\qedsymbol}{\QED}
\newcommand{\x}{\mathbf{x}}
\newcommand{\limsupf}{\limsup_{\x \rightarrow \x_0; \x \in E} f\left(\x\right)}
\newcommand{\supf}{\sup_{\x \in B'\left(\x_0;\delta\right)\cap E} f\left(\x\right)}
\newcommand{\y}{\mathbf{y}}
\newcommand{\convm}{\overset{m}{\longrightarrow}}
\newcommand{\w}{\omega}
\newcommand{\al}{\alpha}
\newcommand{\eps}{\epsilon}
\newcommand{\ph}{\varphi}
\newcommand{\tab}{\hspace*{2em}}
\newcommand{\bigsp}{\vspace{1 in}}
\newcommand{\smallsp}{\vspace{.5 in}}

\begin{document}
\pagestyle{empty}


\begin{center}
\begin{large}
Math 107-01\\
Test 3A: November 4, 2013
\end{large}
\end{center}

\begin{small} Answer all of the following questions.  You will have 50 minutes to complete the
test. You may not use any outside assistance on this test.  You may use a
calculator, but you may not program information into the calculator. The use of
electronic equipment such as mp3 players, ipods, cell phones and any
electronic devices other than your calculator during the test is
prohibited.  If you have time, check your answers for accuracy.  \textbf{Show all work. Unsupported answers may not
receive credit.}  Good luck!\end{small}

\vspace{ .15 in}
Name: \underline{\qquad \qquad \qquad \qquad \qquad \qquad \qquad \qquad \qquad \qquad}\\

\begin{enumerate}
\item A die is rolled and a coin flipped.  Find the following (3 points each):
\begin{enumerate}
\item The sample space
\vspace{.9 in}
\item The number of outcomes in the sample space
\vspace{.7 in}
\end{enumerate}

\item A card is dealt from a standard deck of 52 cards.  Find the probability of being dealt each of the following.  State your answers as reduced fractions (4 points each):
\begin{enumerate}
\item a four
\vspace{.4 in}
\item a diamond
\vspace{.4 in}
\item not a diamond
\vspace{.4 in}
\item the four of diamonds
\vspace{.4 in}
\item a four or a diamond
\vspace{.4 in}

\end{enumerate}

\item A die is rolled.  Find the following (3 points each):
\begin{enumerate}
\item the probability of rolling a 2
\vspace{.5 in}
\item the odds of rolling a 2
\vspace{.5 in}
\item the probability of rolling a number greater than 4
\vspace{.5 in}
\item the odds of rolling a number greater than 4
\vspace{.5 in}
\end{enumerate}



\item
A jar contains 35 jellybeans.  There are twelve blue, eight red, ten yellow, and five green jellybeans.  You pick 5 jellybeans without looking.  Find the probability of each of the following.  Show all work and round your answers to four decimal places (4 points each):

\begin{enumerate}
\item
all five are blue

\vspace{.6 in}



\item
exactly four are blue

\vspace{.6 in}

\item
at least four are blue

\vspace{.4 in}
\end{enumerate}

\newpage
\item (4 points each)
\begin{enumerate}
\item
For a certain factory, the probability of selecting a worker earning a particular hourly wage is as follows:

\begin{tabular}{|c|c|c|c|c|c|c|}
\hline
Hourly wage	& \$8.50 & \$9.00 & \$9.50 & \$10.00 & \$10.50 & \$ 11.00 \\
\hline
Probability & 15\% & 20\% & 25\% & 20\% & 15\% & 5\%\\
\hline
\end{tabular}
\\Find the expected hourly wage that a worker at this factory makes.  Round to the nearest cent.

\vspace{1.1 in}
\item Find the expected value of the following game:
\\You roll a die.  If you roll a 2,4, or 6, then you win \$15.
\\If you roll a 3 or 5, then you lose \$10.
\\If you roll a 1, then you lose \$20.
\\Round to the nearest cent.
\end{enumerate}
\vspace{1.1 in}

\item
In a certain Math 107 class, 75\% of the students live on campus.  80\% of the students are taking an English course.  60\% of the students live on campus and are taking an English course.

Let L=\{live on campus\} and E=\{taking an English course\}.

Construct a Venn diagram illustrating the situation and find $p(L\cup E)$. (5 points)

\includegraphics{venn1test3.jpg}

$p(L\cup E)=\underline{\hspace{1.0 in}}$

\item
$A$ and $B$ are two events for a sample space. $p(A)=0.58$ and $p(B) = 0.29$.  You must show \underline{all} work to receive \underline{full} credit. (3 points each):
\begin{enumerate}
\item If $p(A\cap B)=0.16$, what is $p(A\cup B)$?

\vspace{1.3 in}

\item If $A$ and $B$ are mutually exclusive events, what is $p(A\cup B)$?
\vspace{1.4 in}
\end{enumerate}

\item
The probability that the new library is finished in time for the Summer Festival is 55\%.  If the library is finished for the Summer Festival, the probability it will have a roof that leaks is 25\%.  If the library is not finished in time for the Summer Festival, the probability that it will have a roof that leaks is 10\%.

Construct a tree diagram, including the probability of each limb.  (10 points)

\includegraphics{venn3test3.jpg}

\newpage
\item
Of the students who are Biology majors, 55\% are taking Calculus and 65\% are taking Ecology.  30\% of the Biology majors are taking both courses.

Let C=\{taking Calculus\} and E=\{taking Ecology\}.

The Venn diagram is given below.

\includegraphics{venn2test3.jpg}

Jack is a Biology major.  Calculate the following probabilities (3 points each):
\begin{enumerate}
\item
Jack is not taking Calculus
\vspace{.3 in}
\item
Jack is taking Calculus and Ecology
\vspace{.3 in}
\item
Jack is taking Calculus or Ecology
\vspace{.3in}
\item
Jack is not taking Calculus and not taking Ecology
\vspace{.3 in}
\item
Jack is not taking Calculus and is taking Ecology
\vspace{.3 in}
\item
Jack is taking Calculus given that he is taking Ecology
\vspace{.3 in}
\item
Jack is taking Ecology given that he is taking Calculus
\vspace{.3 in}
\end{enumerate}
\vspace{.5 in}
On my honor, I have neither given nor received any academic aid or information that would violate the Honor Code of Mars Hill University. \underline{\hspace{2.0 in}}(please sign)
\end{enumerate}
\end{document} 